\begin{textsample}{POS Dim 3 – es – Score 20.00 – es/es\_inf\_3\_direitos.txt}  \label{ex:f3_pos_116}
Direitos de aprendizagem Constituem como pontos referenciais os 6 ( seis ) direitos de aprendizagem que os \textbf{bebês}, as \textbf{crianças} bem \textbf{pequenas} e as \textbf{crianças} \textbf{pequenas} devem ter garantidos : - Conviver com outras \textbf{crianças} e \textbf{adultos}, em \textbf{pequenos} e grandes grupos, utilizando diferentes linguagens, ampliando o conhecimento de si e do outro, o respeito em relação à cultura e às diferenças entre as pessoas. - \textbf{Brincar} cotidianamente de diversas formas, em diferentes espaços e tempos, com diferentes \textbf{parceiros} ( \textbf{crianças} e \textbf{adultos} ), ampliando e diversificando seu acesso a produções culturais, seus conhecimentos, sua imaginação, sua criatividade, suas experiências emocionais, corporais, \textbf{sensoriais}, expressivas, cognitivas, sociais e \textbf{relacionais}. - Participar ativamente, com \textbf{adultos} e outras \textbf{crianças}, tanto do planejamento da gestão da escola e das atividades propostas pelo educador quanto da realização das atividades da vida cotidiana, tais como a escolha das \textbf{brincadeiras}, dos materiais e dos ambientes, desenvolvendo diferentes linguagens e elaborando conhecimentos, decidindo e se posicionando. - Explorar movimentos, \textbf{gestos}, \textbf{sons}, formas, \textbf{texturas}, cores, palavras, emoções, transformações, relacionamentos, histórias, objetos, elementos da natureza, na escola e fora dela, ampliando seus saberes sobre a cultura, em suas diversas modalidades : as artes, a escrita, a ciência e a tecnologia. - Expressar, como sujeito dialógico, criativo e sensível, suas necessidades, emoções, sentimentos, dúvidas, hipóteses, descobertas, opiniões, questionamentos, por meio de diferentes linguagens. - Conhecer -se e construir sua identidade pessoal, social e cultural, constituindo uma imagem positiva de si e de seus grupos de pertencimento, nas diversas experiências de \textbf{cuidados}, interações, \textbf{brincadeiras} e linguagens vivenciadas na \textbf{instituição} escolar e em seu contexto familiar e comunitário. ( BNCC, p. 36, 2017 ) As características e peculiaridades das \textbf{crianças}, que hoje frequentam as unidades de Educação \textbf{Infantil}, devem impulsionar os educadores a garantir que os direitos de aprendizagem sejam mediadores de significativas aprendizagens, conferindo intencionalidade às práticas pedagógicas, a fim de contemplar as diferenças e diversidades, características da \textbf{infância}.
Assim, faz -se importante alinhar, com os profissionais de cada \textbf{instituição}, os conceitos sobre a \textbf{infância}, tempos, espaços escolares, sobre as concepções teóricas adotadas para a elaboração de currículos, que respondam aos questionamentos e subsidiem intervenções intencionais nos processos de ensino e aprendizagem, considerando : > Essa concepção de \textbf{criança} como ser que observa, questiona, levanta hipóteses, conclui, faz julgamentos e \textbf{assimila} valores e que constrói conhecimentos e se apropria do conhecimento sistematizado por meio da ação e nas interações com o mundo físico e social não deve resultar no confinamento dessas aprendizagens a um processo de desenvolvimento natural ou espontâneo.
Ao contrário, impõe a necessidade de imprimir intencionalidade educativa às práticas pedagógicas na Educação \textbf{Infantil}, tanto na \textbf{creche} quanto na \textbf{pré-escola}. ( BRASIL, 2017 – p. 34, grifo no original ) A intenção educativa deve consubstanciar tanto o planejamento, a organização do ambiente pedagógico pelo professor e o \textbf{acompanhamento} e avaliação da aprendizagem e desenvolvimento da \textbf{criança}.

% matched lemmas: acompanhamento, adulto, assimilar, bebê, brincadeira, brincar, creche, criança, cuidado, gesto, infantil, infância, instituição, parceiro, pequeno, pré-escola, relacional, sensorial, som, textura
\end{textsample}
